% 图契约：麦克斯韦方程组中两条散度方程与两条旋度方程的物理分工；几何、方向和公式与正文同一约定。
\documentclass[tikz,border=5pt]{standalone}
\usepackage{fontspec,xeCJK,amsmath,amssymb}
\IfFontExistsTF{Times New Roman}{\setmainfont{Times New Roman}}{\setmainfont{TeX Gyre Termes}}
\IfFontExistsTF{Heiti SC}{\setCJKmainfont[AutoFakeBold=1.5]{Heiti SC}}{\setCJKmainfont[AutoFakeBold=1.5]{FandolHei-Regular}}
\usetikzlibrary{arrows.meta,calc,patterns}
\definecolor{ink}{HTML}{243745}
\definecolor{blue}{HTML}{275DAB}
\definecolor{green}{HTML}{177C63}
\definecolor{red}{HTML}{BD4935}
\definecolor{orange}{HTML}{AA6C16}
\definecolor{purple}{HTML}{715393}
\begin{document}
\begin{tikzpicture}[x=1cm,y=1cm,font=\fontsize{9.5}{12}\selectfont,
v/.style={-{Latex[length=2mm,width=1.3mm]},line width=.85pt},
axis/.style={-{Latex[length=1.7mm,width=1.1mm]},ink!65,line width=.55pt},
aux/.style={ink!45,dashed,line width=.55pt},
lab/.style={fill=white,inner sep=1.8pt},
card/.style={draw=ink!25,fill=blue!3,rounded corners=2pt,text width=6.1cm,minimum height=.85cm,inner sep=6pt,align=center}]
\path[use as bounding box](-8,-3.6)rectangle(8,3.65);
\node[font=\bfseries\fontsize{10.5}{13}\selectfont,text width=15.1cm,align=center]at(0.00000,3.22000){M7-H11　麦克斯韦方程组中两条散度方程与两条旋度方程的物理分工};
\begin{scope}[shift={(-4,0)}]
\draw[blue](0,1.15)circle(0.8);
\node[text=red]at(0.00000,1.15000){$q$};
\draw[v,blue](0.85000,1.15000)--(1.30000,1.15000);
\draw[v,blue](0.00000,2.00000)--(0.00000,2.45000);
\draw[v,blue](-0.85000,1.15000)--(-1.30000,1.15000);
\draw[v,blue](-0.00000,0.30000)--(-0.00000,-0.15000);
\draw[green](0,-1.35)circle(0.8);
\draw[v,green](.8,-1.35)arc(0:250:.8);
\draw[red](0,-1.35)circle(0.13);
\fill[red](0,-1.35)circle(.035);
\node[]at(1.90000,1.15000){通量};
\node[]at(1.90000,-1.35000){环量};
\end{scope}
\node[card]at(4.00000,1.25000){$\nabla\cdot\mathbf D=\rho_f,\quad\nabla\cdot\mathbf B=0$};
\node[card]at(4.00000,-1.05000){$\nabla\times\mathbf E=-\dot{\mathbf B},\quad\nabla\times\mathbf H=\mathbf J_f+\dot{\mathbf D}$};
\end{tikzpicture}\end{document}
